\section{METODOLOGÍA Y DISEÑO DE LA INVESTIGACIÓN}

\USFQNotaPlantilla{Esta sección debe contener: explicación de la metodología que se utilizará de acuerdo con las preguntas de investigación planteadas; explicación clara de por qué la metodología elegida responde a la pregunta de investigación; detalle de la herramienta de investigación, por ejemplo encuesta, experimento, entrevistas, documentos o estudio de caso; información sobre participantes si aplica; descripción de cómo se recolectaron los datos; explicación de cómo fueron analizados los datos y los sistemas de criterios y codificación aplicados.}

\USFQNotaPlantilla{Un aspecto clave de esta sección es que deben proveer todas las descripciones necesarias para que el estudio pueda ser replicado de la manera más fidedigna posible, de tal manera que se puedan revisar los resultados sin infringir sesgos metodológicos.}

\subsection{Justificación de la metodología}

\USFQCampo{Explicar la metodología seleccionada y justificar su correspondencia con la pregunta de investigación.}

\subsection{Herramientas de investigación}

\USFQCampo{Describir los instrumentos, procedimientos, experimentos, entrevistas, documentos, protocolos o herramientas utilizadas.}

\subsection{Participantes, muestras o documentos}

\USFQCampo{Si aplica, indicar número de participantes, criterios de selección, variables demográficas relevantes y límites del estudio. Si no aplica, describir la muestra documental, experimental o de laboratorio.}

\subsection{Recolección y análisis de datos}

\USFQCampo{Detallar origen de los datos, características de la muestra, prueba piloto si corresponde, procedimientos de análisis, criterios de codificación y controles de calidad.}
